VisionServe makes the model license policy a first-class part of the serving
infrastructure: every model is admitted through a strict Apache-2.0/MIT/BSD allowlist,
so a deployment cannot silently inherit copyleft obligations. Two further design choices
support this gate. The system ships as a single Go binary with no Python runtime
dependency, and it routes all inference exclusively through ONNX Runtime~\cite{onnxruntime}
behind a unified \texttt{Result} schema shared across all task types. Together these
choices make deployment reproducible on hardware ranging from a Jetson Orin to an
Apple Silicon laptop, without operator intervention.

We evaluated the system across 19 registered models covering seven task types:
object detection, instance segmentation, depth estimation, face detection, OCR,
open-vocabulary detection, and planar grasp synthesis. Measured end-to-end latencies
range from 33\,ms to 570\,ms on an RTX A6000 GPU, and the majority of models stay
under 100\,ms at typical input resolutions. Reproduced on a Jetson AGX Orin, the same
single binary drives the TensorRT/CUDA/CPU fallback at up to 20$\times$ the CPU throughput
and 12$\times$ the energy efficiency for a real detector.

VisionServe is released under the Apache-2.0 license and is publicly available at
\url{https://github.com/mtbui2010/visionserve}. We invite community contributions: a new
model requires only a single package under \texttt{internal/models/} implementing the
\texttt{Model} or \texttt{PipelineModel} interface, with no changes to core server or
engine code. Planned extensions include the CLIP~\cite{radford2021clip} text encoder,
RTMPose, ArcFace, ROCm execution-provider support, and multi-image request batching.
